% SCP10, Definition 5.6, Theorem 5.9 and Observation 6.6.
% Native periodic closures and coherent sums; physical periods at least three.

\begin{theorem}[Physical transport uniform in inserted bonds]%
    \label{thm:peps_inserted_physical_support_transport}
    \lean{TNLean.PEPS.graphPhysicalProductMap_graphInsertedBondNetwork,
        TNLean.PEPS.graphInsertedBondNetwork_mem_range_productSiteMap,
        TNLean.PEPS.graphInsertedBondNetwork_mem_euclideanRange_productSiteMap,
        TNLean.PEPS.exists_graphGIsometricInsertedPhysicalSupportTransport}
    \leanok
    Let $A_v,B_v$ be G-isometric tensors for the same unitary virtual
    representation at each vertex of a finite simple graph. There are positive
    factors $c_{A,v},c_{B,v}$, local physical maps $F_v$, and one isometry $I$
    on the full product of the ranges of the $A_v$ such that, for every
    collection of inserted bond matrices $K=(K_e)_e$,
    \begin{align}
        I\Psi_A(K)=\left(\prod_v\frac{\sqrt{c_{A,v}}}{\sqrt{c_{B,v}}}\right)
        \Psi_B(K).
    \end{align}
    Every inserted state belongs to this same product range. The isometry
    acts there by the product of the $F_v$, independently of $K$.
\end{theorem}
\begin{proof}\leanok
    \uses{thm:peps_g_isometric_physical_transport}
    Choose the local normalized-adjoint maps before choosing $K$.
    Their column Gram identities imply isometry on the product range.
    Expand each inserted contraction over its independent endpoint labels.
    Products of physical maps commute with this sum, and every summand
    belongs to the product range. The local tensor identities give the
    displayed equality with the same factors for every $K$.
\end{proof}

\begin{definition}[Permutation bonds under two-by-two grouping]%
    \label{def:peps_twisted_two_by_two_coordinates}
    \lean{TNLean.PEPS.twoByTwoHorizontalPerm,
        TNLean.PEPS.twoByTwoVerticalPerm,
        TNLean.PEPS.twoByTwoPermutationBoundary,
        TNLean.PEPS.pairedLeftRegularMatrix}
    \leanok
    Group the fine torus into disjoint two-by-two squares. Give each pair
    of horizontal crossing bonds the same permutation $s_h$, and each pair
    of vertical crossing bonds the same permutation $s_v$; leave all
    internal bonds unchanged. The paired regular action is
    $L_g^{(2)}=L_g\otimes L_g$, acting on $(a,b)$ by $(ga,gb)$.
    Horizontal insertions act on the receiving left leg, and vertical
    insertions on the top leg in the native downward orientation.
\end{definition}

\begin{theorem}[Exact twisted geometric grouping]%
    \label{thm:peps_twisted_two_by_two_geometry}
    \lean{TNLean.PEPS.twoByTwoSiteLegs_permutation,
        TNLean.PEPS.torusBondNetwork_perm_eq_twoByTwoTiledSum,
        TNLean.PEPS.torusBondNetwork_perm_eq_twoByTwoBlocked,
        TNLean.PEPS.kitaevDoubleCoordinateEquiv_seam,
        TNLean.PEPS.twoByTwoHorizontalPerm_seam,
        TNLean.PEPS.twoByTwoVerticalPerm_seam,
        TNLean.PEPS.pairedLeftRegularMatrix_apply,
        TNLean.PEPS.pairedLeftRegularMatrix_eq_kronecker,
        TNLean.PEPS.torusGClosure_leftRegular_eq_twoByTwoBlocked}
    \leanok
    For every pair of positive coarse periods $w,h$, the actual fine
    contraction with the crossing permutations above equals the actual
    coarse contraction of the four-site tensors with paired permutations.
    In particular, a native closure $(g_0,h_0)$ becomes precisely the same
    closure $(g_0,h_0)$ for the paired representation. No scalar occurs.
    This geometric identity requires neither symmetry nor commutation of
    the closure labels, and includes periods one and two.
\end{theorem}
\begin{proof}\leanok
    \uses{def:peps_twisted_two_by_two_coordinates}
    Reindex all fine bonds into four internal bonds per block and the two
    labels of every coarse crossing bond. The corner-by-corner endpoint
    identity retains the original heads of both orientations. Summing the
    internal labels gives exactly the four-site tensor. For a doubled
    periodic coordinate $2x+b$, the outgoing bond crosses the seam precisely
    when $b=1$ and $x$ is the final coarse coordinate. Thus precisely the
    two crossing bonds, and no internal bond, receive each seam insertion.
\end{proof}

\begin{definition}[Ordered labels for native torus seams]%
    \label{def:peps_twisted_torus_ordered_labels}
    \lean{TNLean.PEPS.torusGraphRegularLabels,
        TNLean.PEPS.torusClosureRegularLabels}
    \leanok
    For a simple torus graph with periods at least three, express every
    native bond in the ordered-edge orientation. A reversed native arrow
    replaces its group label $u$ by $u^{-1}$. The horizontal periodic seam
    carries $h_0$ and the downward vertical seam carries $g_0$.
\end{definition}

\begin{theorem}[Native, paired, and ordered inserted contractions]%
    \label{thm:peps_twisted_torus_graph_coordinates}
    \lean{TNLean.PEPS.graphInsertedBondNetwork_equiv,
        TNLean.PEPS.graphInsertedBondNetwork_torusIncidentFamily,
        TNLean.PEPS.graphInsertedBondNetwork_twoByTwoBundledGraphSite,
        TNLean.PEPS.torusGraphBondMatrix_leftRegular,
        TNLean.PEPS.torusGraphBondMatrix_regularBundle,
        TNLean.PEPS.graphInsertedBondNetwork_torusClosureRegularLabels,
        TNLean.PEPS.regularBundleMatrix_unit_eq_paired,
        TNLean.PEPS.torusGClosure_leftRegular_eq_twoByTwoBundledGraphSite}
    \leanok
    Native torus contractions equal their ordered-edge inserted graph
    contractions, with arbitrary site-dependent physical alphabets.
    Bijective relabeling of both endpoints preserves every coefficient.
    Hence the native fine closure equals the inserted coarse contraction
    of its actual four-site tensors, with two regular labels on each bond.
    The same ordered labels describe a single regular bond and a bundle
    of regular bonds: matrix transposition under reversal is exactly group
    inversion in both cases.
\end{theorem}
\begin{proof}\leanok
    \uses{def:peps_twisted_torus_ordered_labels,
        thm:peps_twisted_two_by_two_geometry}
    Use the incidence-coordinate and endpoint bijections in the finite
    sums. A permutation matrix for left multiplication satisfies
    $L_u^{\mathsf T}=L_{u^{-1}}$, also for simultaneous multiplication of
    the two labels. The resulting native and ordered contraction factors
    therefore agree at every endpoint, including periodic seam endpoints.
\end{proof}

\begin{theorem}[Inserted bonds preserve the surplus Bell state]%
    \label{thm:peps_inserted_regular_surplus}
    \lean{TNLean.PEPS.regularBundleMatrix_transpose,
        TNLean.PEPS.regularSurplusCoordinates_twisted,
        TNLean.PEPS.graphInsertedBondNetwork_regularGraphSurplusSite,
        TNLean.PEPS.regularGraphOriginalOutputRegrouping,
        TNLean.PEPS.exists_regularGraphInsertedOriginalTensorIsometry}
    \leanok
    On a bundle $(a,(b_k)_k)$, pass to relative coordinates
    $(a,(a^{-1}b_k)_k)$. An arbitrary oriented insertion $u_e$ acts only
    on the distinguished coordinate. Thus the inserted contraction of the
    original tensor with surplus identity registers equals its original
    inserted contraction times the untwisted surplus equality vectors.
    A single physical-support isometry transports every bundled inserted
    state to this form, with the same local normalization factors and the
    same normalized Bell product.
\end{theorem}
\begin{proof}\leanok
    \uses{thm:peps_inserted_physical_support_transport,
        thm:peps_regular_graph_surplus_factorization}
    At a twisted head the relative label remains unchanged because
    $(ua)^{-1}(ub_k)=a^{-1}b_k$. The sum over distinguished labels gives
    the original inserted contraction, and the independent relative sums
    give the surplus equality vectors. Apply uniform physical support
    transport and retain their inverse-square-root normalizations.
\end{proof}

\begin{theorem}[Nonzero native closures in one fine physical support]%
    \label{thm:peps_twisted_two_by_two_support}
    \lean{TNLean.PEPS.IsGIsometric.torusLegRep_of_regularFourLeg,
        TNLean.PEPS.IsGIsometric.regularFourLegClosure_ne_zero,
        TNLean.PEPS.twoByTwoFineClosure_mem_physicalSupport,
        TNLean.PEPS.twoByTwoSupportedClosure}
    \leanok
    Let $A$ be a regular G-isometric four-leg tensor. Every actual native
    closure is nonzero. On a doubled torus with coarse periods at least
    three, all fine closure vectors belong to the same product block
    physical support, pulled back by the full physical grouping unitary.
\end{theorem}
\begin{proof}\leanok
    \uses{thm:peps_twisted_torus_graph_coordinates,
        thm:peps_inserted_physical_support_transport}
    Identify the four-coordinate tensor with its native four-leg form.
    Its closure self-overlap is a positive site factor times a positive
    power of $|G|$ and the order of the common centralizer of the closure
    labels. This centralizer contains the identity. Support membership
    follows from the exact inserted coarse contraction and its product
    range membership.
\end{proof}

\begin{theorem}[Uniform twisted fixed point and coherent superpositions]%
    \label{thm:peps_twisted_two_by_two_fixed_point}
    \lean{TNLean.PEPS.physicalStateProduct_sum,
        TNLean.PEPS.LinearIsometry.coherent_physicalStateProduct,
        TNLean.PEPS.exists_regularTwoByTwoTwistedTensorIsometry}
    \leanok
    Let $A$ be a homogeneous regular G-isometric tensor and let both coarse
    periods be at least three. Write $\Psi_f(g,h)$ and $\Psi_c(g,h)$ for its
    actual fine and coarse native closure vectors. There is one isometry
    $I$ on the full fine block physical support, acting by physical grouping
    and fixed local physical maps, such that simultaneously for every $g,h$,
    \begin{align}
        I\Psi_f(g,h) & =R\,\Psi_c(g,h)\otimes\Omega,
                     & I\widehat{\Psi_f(g,h)}        & =\widehat{\Psi_c(g,h)}\otimes\Omega,
    \end{align}
    where $R>0$ retains all block and Bell factors and $\|\Omega\|=1$.
    The coarse tensor is literally $A$. For arbitrary complex
    coefficients $z_{g,h}$, put $x=\sum_{g,h}z_{g,h}\Psi_f(g,h)$ and
    $y=\sum_{g,h}z_{g,h}\Psi_c(g,h)$. Then
    \begin{align}
        Ix=R\,y\otimes\Omega,\qquad
        x\ne0\ \Longleftrightarrow\ y\ne0,\qquad
        I\widehat x=\widehat y\otimes\Omega
    \end{align}
    whenever these vectors are nonzero. In particular this holds for
    arbitrary coherent sums supported on commuting closure pairs, retaining
    all their relative coefficients and phases.
\end{theorem}
\begin{proof}\leanok
    \uses{thm:peps_twisted_two_by_two_support,
        thm:peps_inserted_regular_surplus,
        thm:peps_two_by_two_original_scale,
        thm:peps_physical_state_normalization}
    Apply the uniform inserted support isometry between the actual
    four-site block tensor and the original tensor with surplus registers.
    Exact native seam coordinates identify its input and output. Their
    positive norm relation gives the normalized equalities. Linearity
    gives the coherent equality with the same $I$ and $R$; its norm relation
    is $\|x\|=R\|y\|$, proving nonvanishing equivalence and normalization.
\end{proof}

The all-pairs algebraic identity asserts a ground-state interpretation only
for commuting closures. This is the native periodic, physical-support
specialization of \cite[Observation~6.6]{Schuch2010PEPS}. The geometric
identity includes all positive periods, but physical Bell separation for
coarse periods one and two, other boundary conditions, and extension to an
ambient physical unitary remain outside this statement.

% Remaining-scope note: docs/paper-gaps/scp10_twisted_two_by_two_regular_fixed_point.tex.

The physical map is chosen before the seam labels and coefficients.
The algebraic identity includes noncommuting labels, while the
parent-ground-state interpretation uses commuting closures. The remaining
period and boundary restrictions are recorded in
\cite{gap:scp10_twisted_two_by_two_regular_fixed_point}.
